% !TeX program = pdflatex
% C(13) >= 36 for the "no 5 on a sphere" grid problem (AlphaEvolve Problem 6.60)
% Build: pdflatex note.tex (run twice). See BUILD.md.
\documentclass[10pt]{article}

\usepackage[T1]{fontenc}
\usepackage{lmodern}
\usepackage[margin=1in]{geometry}
\usepackage{amsmath,amssymb,amsthm}
\usepackage{booktabs}
\usepackage{microtype}
\usepackage[hidelinks]{hyperref}
\usepackage{xurl}

\newtheorem{theorem}{Theorem}[section]
\newtheorem{proposition}[theorem]{Proposition}
\newtheorem{lemma}[theorem]{Lemma}
\theoremstyle{remark}
\newtheorem{remark}[theorem]{Remark}

\newcommand{\Z}{\mathbb{Z}}
\newcommand{\R}{\mathbb{R}}

\title{A 36-point subset of the $13\times 13\times 13$ grid\\
with no five points on a sphere or a plane}
\author{John Erlbacher\thanks{Independent Researcher. Email: \texttt{erlbacher.research@gmail.com}. ORCID: 0009-0003-6851-4139.}}
\date{June 2026}

\begin{document}
\maketitle

\begin{abstract}
For the ``no 5 on a sphere'' problem --- determine $C(n)$, the size of the
largest subset of the grid $\{1,\dots,n\}^3$ containing no five points on a
common sphere or plane --- the strongest published lower bounds for
$7 \le n \le 12$ are due to AlphaEvolve, ending with $C(12) \ge 33$; nothing
specific to $n = 13$ has been published, so the best previously known lower
bound was $C(13) \ge C(12) \ge 33$. We exhibit an explicit $36$-point subset of
$\{0,\dots,12\}^3$ with no five points on a sphere or a plane, proving
$C(13) \ge 36$. The set is invariant under the central inversion
$p \mapsto (12,12,12) - p$: eighteen antipodal point pairs whose squared
distances from the cube center are pairwise distinct. Verification is a
finite exact-integer computation --- all $\binom{36}{5} = 376{,}992$ lifted
$5 \times 5$ determinants are nonzero --- reproduced here by a $21$-line
program, printed in full, and carried out by three independently written
checkers validated by mutation testing. The set was found in $841$ seconds
on eight CPU threads by an exact-arithmetic local search restricted to
centrally symmetric configurations; en route, four independent search
pilots produced thousands of distinct $34$-point sets and dozens of
$35$-point sets. The trivial upper bound $C(13) \le 52$ is unaffected.
\end{abstract}

\section{Introduction}\label{sec:intro}

The following problem appears as Problem~6.60 in the AlphaEvolve
``repository of problems'' of Georgiev, G\'omez-Serrano, Tao and Wagner
\cite{AlphaEvolveMath,Repo}:

\begin{quote}
\emph{``For $n$ a natural number, let $C(n)$ denote the size of the largest
subset of $[n]^3 = \{1,\dots,n\}^3$ such that no 5 points lie on a sphere or
a plane. Obtain upper and lower bounds for $C(n)$ that are as strong as
possible.''}
\end{quote}

\noindent
It is the three-dimensional case of the no-$(d{+}2)$-on-a-sphere problem
(see \cite[Problem~4 of Chapter~10.1]{BMP2005}), a generalization, going
back to Thiele \cite{Thiele1995}, of the no-four-on-a-circle problem of
Erd\H{o}s and Purdy. Thiele \cite{Thiele1995} proved
$c\sqrt{n} \le C(n) \le 4n$ (the upper bound is the observation that each of
the $n$ planes $x = \text{const}$ carries at most $4$ points of a valid
set); Suk and White \cite{SukWhite2025} improved the lower bound to
$n^{3/4-o(1)}$, and Dong and Xu \cite{DongXu2025} to $n - o(n)$, with a
related construction from rational normal curves given by Szab\'o
\cite{Szabo2025}. These constructions are asymptotic and yield no values for
specific small $n$.

For small $n$ the record lower bounds come from AI-driven search. PatternBoost
\cite{PatternBoost} reported $C(3) \ge 8$, $C(4) \ge 11$, $C(5) \ge 14$,
$C(6) \ge 18$, $C(7) \ge 20$, $C(8) \ge 22$, $C(9) \ge 25$ and
$C(10) \ge 27$;\footnote{The explicit witnesses printed in the appendix of
\cite{PatternBoost} pass the determinant criterion of
Proposition~\ref{prop:det} for $n = 3,\dots,9$ (we re-checked them), but the
printed $n = 10$ list is corrupted --- as printed it contains coordinates up
to $13$ and a vanishing $5$-subset --- so we quote that value as reported; it
is in any case superseded by AlphaEvolve's verified $C(10) \ge 28$.}
AlphaEvolve \cite{AlphaEvolve,AlphaEvolveMath} improved
these to $C(7) \ge 21$, $C(8) \ge 23$, $C(9) \ge 26$, $C(10) \ge 28$, and
added the new bounds $C(11) \ge 31$ and $C(12) \ge 33$
\cite[Problem~6.60]{AlphaEvolveMath}. Beyond the general bounds above and the
monotone consequence $C(13) \ge C(12) \ge 33$, we could find no published
bound for any specific $n \ge 13$ (literature and web searches dated 11--12
June 2026; see Section~\ref{sec:scope}), so the baseline for $n = 13$ was
$33$. This note establishes:

\begin{theorem}\label{thm:main}
$C(13) \ge 36$. Explicitly, the $36$-point set $S \subseteq \{0,\dots,12\}^3$
of Table~\ref{tab:cert} contains no five points on a common sphere or plane.
\end{theorem}

By monotonicity, $C(n) \ge 36$ for all $n \ge 13$. The best known upper
bound remains the trivial $C(13) \le 4 \cdot 13 = 52$.

\paragraph{Conventions.}
We work in $\{0,\dots,12\}^3$; adding $1$ to every coordinate places $S$ in
the problem's $[13]^3$, and the defining condition is translation-invariant.
``No five on a sphere or a plane'' is interpreted, as in the source problem
and its official evaluation code \cite{Repo}, through the determinant
criterion of Proposition~\ref{prop:det}, which treats planes as degenerate
spheres; the checkers used here accept all six published AlphaEvolve record
sets for $n = 7,\dots,12$ under this criterion
(Section~\ref{sec:verification}), so the comparison $33 \to 36$ is made
under the very convention that produced the $33$.

\paragraph{Methods and AI-use disclosure.}
The set was found by Demonstrandum, a verification-first multi-agent AI
search pipeline built on Claude language models (Anthropic): agents froze the problem statement and
conventions from the official sources, proposed attack plans, and wrote and
ran the search programs (Rust and Python) of Section~\ref{sec:method}; every
geometric test in the search and every step of the verification is exact
integer arithmetic (floating point appears only in search-heuristic branch
probabilities and timers, never in any validity decision). Certificates
were verified by three independently written exact-integer checkers,
validated by mutation testing, plus a fourth written by an adversarial
referee agent from a different model family (Codex, GPT-5.5, OpenAI)
(Section~\ref{sec:verification}). The text of this note was likewise
drafted with the Claude-based pipeline; the author directed the work and
takes responsibility for its content. Every number in this note that derives
from the certificate or from our own runs was recomputed from the underlying
artifacts during its preparation (12 June 2026). The mathematical claim of Theorem~\ref{thm:main} rests
solely on the explicit set of Table~\ref{tab:cert} and the finite
computation of Section~\ref{sec:verification}, which any reader can
reproduce in under a minute.

\section{The certificate}\label{sec:cert}

Throughout, identify a point $p = (x,y,z) \in \R^3$ with its \emph{lift}
$\Lambda(p) = (x,\, y,\, z,\, x^2{+}y^2{+}z^2,\, 1) \in \R^5$.

\begin{proposition}\label{prop:det}
Five pairwise distinct points $p_1,\dots,p_5 \in \R^3$ lie on a common
sphere or plane if and only if
$\det\bigl(\Lambda(p_1),\dots,\Lambda(p_5)\bigr) = 0$.
\end{proposition}

\begin{proof}
A sphere or plane is the zero set of
$f(x,y,z) = a(x^2{+}y^2{+}z^2) + bx + cy + dz + e$ for some
$(a,b,c,d,e) \neq 0$: a plane when $a = 0$ (then $(b,c,d) \neq 0$, else
$f \equiv e \ne 0$ has empty zero set), a sphere of radius
$r$, $r^2 = (b^2{+}c^2{+}d^2{-}4ae)/4a^2$, when $a \neq 0$. If the five
points lie on such a zero set, the vector $(b,c,d,a,e)$ is orthogonal to all
five lifts, so the determinant vanishes. Conversely, if the determinant
vanishes, some $(b,c,d,a,e) \neq 0$ is orthogonal to all five lifts, giving
an $f$ as above vanishing at all five points; if $a \ne 0$, its zero set
contains five distinct points, hence is neither empty nor a single point,
hence is a genuine sphere ($r^2 > 0$); if $a = 0$ it is a plane.
\end{proof}

Call a finite set $S \subset \R^3$ \emph{valid} if every $5$-subset has
nonvanishing lifted determinant. By Proposition~\ref{prop:det} this is
exactly the problem's condition. One consequence used heavily by the search
(but not needed for verification) is worth recording:

\begin{lemma}\label{lem:degenerate}
If four distinct points are collinear, or lie on a common circle, then these
four points together with \emph{any} fifth point have vanishing lifted
determinant.
\end{lemma}

\begin{proof}
A line is contained in two distinct planes, and a circle in a plane and a
sphere; in both cases the four lifts satisfy two linearly independent
orthogonality constraints, so they span a subspace of dimension at most $3$.
Together with the lift of any fifth point they therefore lie in a subspace
of dimension at most $4$, and the $5 \times 5$ determinant vanishes.
\end{proof}

Thus a valid set with at least $5$ points has at most $3$ points on any line
and at most $3$ on any circle, and a search must treat such
\emph{rank-degenerate} quadruples as blocking every cell.

\begin{table}[ht]
\centering\small
\begin{tabular}{@{}rllr@{\hspace{2.5em}}rllr@{}}
\toprule
\# & $p$ & $\sigma(p)$ & $4\rho^2$ & \# & $p$ & $\sigma(p)$ & $4\rho^2$ \\
\midrule
1 & $(3,5,7)$ & $(9,7,5)$ & 44 & 10 & $(0,3,7)$ & $(12,9,5)$ & 184 \\
2 & $(4,3,5)$ & $(8,9,7)$ & 56 & 11 & $(3,1,10)$ & $(9,11,2)$ & 200 \\
3 & $(4,8,3)$ & $(8,4,9)$ & 68 & 12 & $(2,6,12)$ & $(10,6,0)$ & 208 \\
4 & $(3,7,10)$ & $(9,5,2)$ & 104 & 13 & $(1,4,1)$ & $(11,8,11)$ & 216 \\
5 & $(1,6,4)$ & $(11,6,8)$ & 116 & 14 & $(0,9,10)$ & $(12,3,2)$ & 244 \\
6 & $(5,2,2)$ & $(7,10,10)$ & 132 & 15 & $(5,0,0)$ & $(7,12,12)$ & 292 \\
7 & $(5,11,3)$ & $(7,1,9)$ & 140 & 16 & $(1,12,11)$ & $(11,0,1)$ & 344 \\
8 & $(6,0,4)$ & $(6,12,8)$ & 160 & 17 & $(0,1,12)$ & $(12,11,0)$ & 388 \\
9 & $(2,8,1)$ & $(10,4,11)$ & 180 & 18 & $(0,12,0)$ & $(12,0,12)$ & 432 \\
\bottomrule
\end{tabular}
\caption{The certificate: a valid $36$-point set
$S \subset \{0,\dots,12\}^3$, listed as the $18$ orbit pairs
$\{p,\sigma(p)\}$ of the central inversion $\sigma(p) = (12,12,12) - p$,
sorted by the squared diameter $4\rho^2 = 4\,|p - c|^2 \in \Z$ of the pair
about the cube center $c = (6,6,6)$. The $18$ values $4\rho^2$ are pairwise
distinct (Lemma~\ref{lem:sym}\,(ii) explains why they must be).}
\label{tab:cert}
\end{table}

The certificate set $S$ is reproduced in machine-readable form here:

{\scriptsize
\begin{verbatim}
[[0,1,12],[0,3,7],[0,9,10],[0,12,0],[1,4,1],[1,6,4],[1,12,11],[2,6,12],
 [2,8,1],[3,1,10],[3,5,7],[3,7,10],[4,3,5],[4,8,3],[5,0,0],[5,2,2],
 [5,11,3],[6,0,4],[6,12,8],[7,1,9],[7,10,10],[7,12,12],[8,4,9],[8,9,7],
 [9,5,2],[9,7,5],[9,11,2],[10,4,11],[10,6,0],[11,0,1],[11,6,8],[11,8,11],
 [12,0,12],[12,3,2],[12,9,5],[12,11,0]]
\end{verbatim}
}

\paragraph{Structure.}
$S$ is invariant under the central inversion $\sigma(p) = (12,12,12) - p$ of
the $13$-cube, with $18$ orbit pairs and no fixed point (the unique fixed
cell $(6,6,6)$ is not in $S$ --- by Lemma~\ref{lem:sym}\,(i) it could not
be). The per-axis layer profiles (number of points with the given
coordinate value $0,1,\dots,12$), palindromic by symmetry, are
\begin{center}\small
$x$: 4,3,2,3,2,3,2,3,2,3,2,3,4;\qquad
$y$: 4,3,1,3,3,2,4,2,3,3,1,3,4;\qquad
$z$: 4,3,4,2,2,3,0,3,2,2,4,3,4;
\end{center}
\noindent
so every axis-aligned plane carries at most $4$ points --- the budget that
the trivial bound $C(13) \le 52$ counts. $16$ of the $36$ points lie on the
surface of the cube. Over all $376{,}992$ lifted $5\times 5$ determinants
the minimum absolute value is $2$ ($200$ subsets have $|\det| \le 10$; the
maximum is $198{,}750$): like the published record sets for $n \le 12$, the
object sits at the very edge of degeneracy. Finally, $S$ is \emph{saturated}:
an exhaustive check of all $2197 - 36$ remaining cells of $\{0,\dots,12\}^3$
(re-run while preparing this note) shows that no $37$th point can be added.

The following elementary facts drove the choice of search space; they are
stated for the $13$-cube but hold for any odd $n$.

\begin{lemma}\label{lem:sym}
Let $c = (6,6,6)$ and $\sigma(p) = 2c - p$, and let $T$ be a
$\sigma$-invariant subset of $\{0,\dots,12\}^3$.
\begin{itemize}\itemsep1pt
\item[(i)] If $c \in T$ and $T$ contains at least two orbit pairs, then $T$
is not valid.
\item[(ii)] If two orbit pairs $\{p, \sigma(p)\}$, $\{q, \sigma(q)\}$ of $T$
satisfy $|p - c| = |q - c|$, then these four points are cocircular or
collinear, so (by Lemma~\ref{lem:degenerate}) $T$ is not valid if
$|T| \ge 5$.
\item[(iii)] Any two orbit pairs of $T$ form a parallelogram, hence a
coplanar quadruple.
\end{itemize}
\end{lemma}

\begin{proof}
(iii): the diagonals $[p, \sigma(p)]$ and $[q, \sigma(q)]$ bisect each other
at $c$. (ii): a parallelogram with equal diagonals is a rectangle (or a
degenerate one: four collinear points), and a rectangle is cyclic. (i): let
$\{p,\sigma(p)\}$ and $\{q,\sigma(q)\}$ be two orbit pairs of $T$. If $p$,
$q$ and $c$ are collinear, all five points lie on one line, hence on one
plane. Otherwise the plane through $c$, $p$ and $q$ is $\sigma$-invariant
(it contains $c$), so it also contains $\sigma(p)$ and $\sigma(q)$: again
five points of $T$ on one plane.
\end{proof}

\section{Verification}\label{sec:verification}

Theorem~\ref{thm:main} reduces, via Proposition~\ref{prop:det}, to a finite
integer computation: \emph{all $\binom{36}{5} = 376{,}992$ five-subsets of
$S$ have nonzero lifted determinant}. Subtracting the first row of the
$5\times 5$ matrix from the others and expanding along the last column
reduces this to a $4\times 4$ integer determinant in the differences
$\Lambda'(p_i) - \Lambda'(p_1)$, $\Lambda'(p) =
(x,y,z,x^2{+}y^2{+}z^2)$. Coordinate differences are at most $12$ and lifted
fourth-coordinate differences at most $432$ in absolute value, so every such
determinant is bounded by $24 \cdot 12^3 \cdot 432 = 17{,}915{,}904 <
2^{25}$: native 64-bit (even 32-bit) integer arithmetic is exact, and no
floating point is needed or used. The complete accept path of the reference
checker is the following $21$ lines of Python (quoted verbatim from
\texttt{check\_cert.py} in the certificate bundle):

{\footnotesize
\begin{verbatim}
import sys, json
from itertools import combinations

def check(points, n=13):
    """Return (True, None) if valid, else (False, reason-or-offending-5-subset)."""
    pts = [tuple(p) for p in points]
    if len(set(pts)) != len(pts):
        return False, "duplicate points"
    if any(len(p) != 3 or any(not isinstance(c, int) or not 0 <= c < n for c in p) for p in pts):
        return False, "non-integer or out-of-range coordinate"
    L = [(x, y, z, x * x + y * y + z * z) for (x, y, z) in pts]
    for idx in combinations(range(len(pts)), 5):
        p = L[idx[0]]
        (a0, a1, a2, a3), (b0, b1, b2, b3), (c0, c1, c2, c3), (d0, d1, d2, d3) = \
            [tuple(L[i][k] - p[k] for k in range(4)) for i in idx[1:]]
        det = ((a0*b1 - a1*b0) * (c2*d3 - c3*d2) - (a0*b2 - a2*b0) * (c1*d3 - c3*d1)
             + (a0*b3 - a3*b0) * (c1*d2 - c2*d1) + (a1*b2 - a2*b1) * (c0*d3 - c3*d0)
             - (a1*b3 - a3*b1) * (c0*d2 - c2*d0) + (a2*b3 - a3*b2) * (c0*d1 - c1*d0))
        if det == 0:
            return False, [pts[i] for i in idx]
    return True, None
\end{verbatim}
}

\noindent
Running \texttt{check\_cert.py} on the list above prints
\texttt{VALID m=36 n=13} in about a second. A skeptical reader is encouraged to re-implement the criterion
independently rather than trust any of our code; the float-based verifier
in the official problem notebook \cite{Repo} should \emph{not} be used as a
certificate checker.

Beyond the reference checker (route A), the certificate bundle contains two
independently written checkers: route B performs fraction-free Bareiss
elimination on the full $5\times 5$ lifted matrix (Python), and route C is a
separate Rust implementation (with $28$ unit tests, including traps built
from the rank-degenerate quadruples of Lemma~\ref{lem:degenerate}). All
three routes were written by different agents than the one that found the
set, and a fourth checker was written from scratch by an adversarial referee
agent from a different model family (Codex, GPT-5.5), which confirmed all
$376{,}992$ determinants nonzero, minimum $|\det| = 2$, and the central
symmetry. The checkers were validated by mutation testing: seven targeted
corruptions of the certificate (a duplicated point; an out-of-range
coordinate $13$; a negative coordinate; a float coordinate; appending the
center cell $(6,6,6)$; replacing a point by the center; appending an
in-range cell completing a coplanar $5$-subset) are rejected by every
applicable route: the float corruption is rejected by both Python routes and
does not even parse for the integer-only Rust route; the other six
corruptions are rejected by all three routes. As a convention anchor, the
same checkers accept all six official
AlphaEvolve record sets for $n = 7,\dots,12$ (sizes $21, 23, 26, 28, 31,
33$, taken verbatim from the official notebook \cite{Repo}) and reject a
known-invalid $34$-point set. All of these checks, and every
certificate- and run-derived number quoted
in this note, were re-run on 12 June 2026 while preparing it.

\section{Method: search in the centrally symmetric subspace}\label{sec:method}

The published $33$-point record set for $n = 12$ is the residue of a generic
evolved local search \cite[Problem~6.60]{AlphaEvolveMath}, with no visible
structure. Our route to $36$ was to \emph{impose} structure: the cube
$\{0,\dots,12\}^3$ has a true center $c = (6,6,6)$, and we searched only
sets invariant under the central inversion $\sigma(p) = 2c - p$, i.e.\
unions of orbit pairs $\{p, \sigma(p)\}$. This halves the search dimension
($36$ points are $18$ pair choices); by Lemma~\ref{lem:sym}\,(iii) every two
orbit pairs automatically form a \emph{coplanar, legal} quadruple,
pre-paying densely the four-points-per-plane budget that record-grade sets
must exhaust, while Lemma~\ref{lem:sym}\,(ii) gives a cheap exact pruning
rule (one squared norm per pair) against the rectangle degeneracies
symmetry would otherwise mass-produce.

The searcher (Rust; exact integer arithmetic in every geometric predicate)
maintains an
incrementally updated \emph{blocker grid}: for the current set $T$ and every
cell $u$ of the grid, the number of $4$-subsets $Q \subseteq T$ whose lifted
cofactor vector is orthogonal to $\Lambda(u)$; by
Proposition~\ref{prop:det}, $u$ is addable only if this count is zero. For a
candidate \emph{pair} $(p, \sigma(p))$, zero blocker counts at both cells
decide all $5$-subsets with at most one new point, and two further exact
layers (cofactor tests over member $2$-subsets, determinant tests over
member $3$-subsets) decide those containing both new points, including all
rank-degenerate cases of Lemma~\ref{lem:degenerate}. On top of this sits a
standard iterated local search: greedy randomized fill over pair orbits,
ruin-and-rebuild (remove $1$--$3$ random pairs, rebuild with a tabu list,
accept if not worse), restart after $800$ non-improving iterations. Every
accepted set is re-verified in-process by the brute-force all-subsets test
before being written to disk.

\begin{table}[ht]
\centering\footnotesize
\begin{tabular}{@{}llll@{}}
\toprule
bound & date found & pilot & evidence \\
\midrule
$C(12) \ge 33$ & Nov.\ 2025 & AlphaEvolve (prior record) & \cite{AlphaEvolveMath,Repo} \\
$C(13) \ge 34$ & 11 June 2026 & centrally symmetric pilot, $4.0$\,s into & run logs \\
 & & its first smoke run ($2$ threads); independently & \\
 & & baseline ILS, $53.8$\,s into its own smoke run; & \\
 & & both further pilots followed within hours & \\
$C(13) \ge 35$ & 11 June 2026 & baseline ILS main run ($t = 538$\,s); & run logs \\
 & & independently, blocker-repair pilot & \\
$C(13) \ge 36$ & 12 June 2026 & centrally symmetric main run: $8$ threads, & Table~\ref{tab:cert} \\
 & & $840.9$\,s elapsed, $611{,}969$ ILS iterations, & \\
 & & $627$ restarts & \\
\bottomrule
\end{tabular}
\caption{Record progression. All searches ran on one desktop CPU (Ryzen~9
9950X3D, 16 cores); roughly three hours elapsed between freezing the
problem dossier and finding the $36$-point set.}
\label{tab:timeline}
\end{table}

Four independent pilots ran in parallel
(Tables~\ref{tab:timeline} and~\ref{tab:abundance}): the
centrally symmetric search; a faithful port of the published evolved
$n = 12$ recipe (the official notebook discloses that, due to a misplaced
\texttt{\#\,EVOLVE-BLOCK-START} line, AlphaEvolve's search for this problem
``had to start its search from scratch every time'' \cite{Repo} --- our port
runs it both as published and with that harness bug fixed); an exact
blocker-cofactor \emph{repair} search over hitting sets of blocking
quadruples; and a finite-field pilot mining mod-$13$ arc structure for
seeds (a direct mod-$13$ construction of size $34$ is impossible: arcs in
$\mathrm{PG}(4,13)$ have at most $14$ points by Ball's theorem
\cite{Ball2012}, so the lifted-matroid-mod-$13$ condition can serve only as
seed diversity).

\begin{table}[ht]
\centering\footnotesize
\begin{tabular}{@{}lllll@{}}
\toprule
pilot & run length & distinct $34$s & distinct $35$s & $36$ \\
\midrule
baseline ILS ($8$ threads, two arms) & $22$\,min snapshot & $198$ & $1$ & --- \\
blocker-repair, $3$ runs & $55 + 48 + 48$\,min & $4068\,/\,2528\,/\,2833$ & $20\,/\,30\,/\,47$ & --- \\
centrally symmetric (symmetric sets) & $14$\,min & $4126$ & n/a & $\mathbf{1}$ \\
finite-field $\mathbb{F}_{13}$, $2$ seeded runs & $5.5 + 11.8$\,min & $69\,/\,71$ & --- & --- \\
\bottomrule
\end{tabular}
\caption{Distinct valid sets logged in the pilots' principal runs (each set
brute-force verified in-process before logging; counts deduplicated within
each run by hash, not across runs and not up to cube symmetries). Of the
baseline ILS $34$s, $182$ of $198$ came from the arm with the harness bug
fixed. A centrally symmetric set has even size, so the symmetric pilot
cannot log a $35$; its $34$s are symmetric ($17$ pairs). The two
finite-field runs' $34$-families overlapped almost completely: all $69$ of
the first run's sets recur among the second's $71$.}
\label{tab:abundance}
\end{table}

Two findings frame the result. First, \emph{abundance}
(Table~\ref{tab:abundance}): valid $34$-point sets in the $13$-cube are not
rare --- pilots logged thousands of distinct ones within minutes, including
$4126$ distinct \emph{centrally symmetric} $34$-sets in the $14$ minutes of
the main symmetric run. Second, \emph{saturation of the old record}: the
published $33$-point set admits no $34$th point in the $13$-cube --- an
exhaustive check of all eight translates of the $12$-cube inside
$\{0,\dots,12\}^3$, re-run while preparing this note, finds zero addable
cells for every translate. The published record was a dead local optimum;
the $34$--$36$ region was simply unexplored. Within the symmetric subspace
the climb did not stop at $34$: where the asymmetric pilots plateaued hard
at $35$ (the first blocker-repair run logged $510$ dead-end local optima at
size $34$ and $18$ at size $35$ in $55$ minutes), one symmetric run pushed
through to $18$ pairs, with the fixed seed $20260612$, reproducible from
the archived search-code snapshot and command line in the bundle.

\section{What this does and does not show}\label{sec:scope}

\begin{itemize}\itemsep0pt\parskip0pt
\item \emph{The exact value of $C(13)$ remains unknown.} Theorem~\ref{thm:main}
is a lower bound; the best upper bound remains the trivial $C(13) \le 52$.
We make no optimality claim.
\item \emph{No empirical maximality claim either.} Our searches are
heuristic. The specific set $S$ admits no $37$th grid point, but other
$36$-point sets, or larger sets, may exist in unexplored basins; no
exhaustive search of any kind is claimed.
\item \emph{Nothing is claimed for other $n$.} In particular we did not find
$C(12) \ge 34$, and we claim no bounds for $n = 14, 15$ beyond the monotone
$C(n) \ge 36$.
\item \emph{No structural theorem.} Central symmetry is a search heuristic
that worked, not a proven property of optimal sets.
\item \emph{Priority.} Openness was checked on 11--12 June 2026 against the
official repository and paper \cite{Repo,AlphaEvolveMath} (whose published
data end at $C(12) \ge 33$), arXiv listings and API, GitHub code search,
OpenReview, and general web search, including an independent sweep by the
adversarial referee agent; no public claim of $C(13) \ge 34$ (or
$C(12) \ge 34$) was found. Concurrent unpublished work cannot be excluded.
\end{itemize}

\paragraph{Data availability.}
The certificate (JSON), the three checkers and the referee's independent
checker, the mutation-test suite, the search code (Rust) with run logs and
the discovering run's command line, and full verification logs will be made
available at \url{https://github.com/demonstrandum-research/[REPO]}. The certificate file
has SHA-256 checksum
\begin{center}
\texttt{\small 333d36ece36e3d845cd2f5bb26e5460f78d881c00418c92cae8bd5215ab0629a}.
\end{center}

\begingroup\small
\begin{thebibliography}{10}

\bibitem{Ball2012}
S.~Ball,
\emph{On sets of vectors of a finite vector space in which every subset of
basis size is a basis},
J. Eur. Math. Soc. \textbf{14} (2012), 733--748.

\bibitem{BMP2005}
P.~Brass, W.~Moser, and J.~Pach,
\emph{Research Problems in Discrete Geometry},
Springer, New York, 2005.

\bibitem{PatternBoost}
F.~Charton, J.~S. Ellenberg, A.~Z. Wagner, and G.~Williamson,
\emph{PatternBoost: constructions in mathematics with a little help from AI},
arXiv:2411.00566 (2024).

\bibitem{DongXu2025}
Z.~Dong and Z.~Xu,
\emph{Large grid subsets without many cospherical points},
arXiv:2506.18113 (2025).

\bibitem{AlphaEvolveMath}
B.~Georgiev, J.~G\'omez-Serrano, T.~Tao, and A.~Z. Wagner,
\emph{Mathematical exploration and discovery at scale},
arXiv:2511.02864 (2025); v3 of 22 December 2025.

\bibitem{Repo}
B.~Georgiev, J.~G\'omez-Serrano, T.~Tao, and A.~Z. Wagner,
\emph{AlphaEvolve repository of problems}, problem~60 (``The no 5 on a
sphere problem''), with evaluation notebook
\texttt{no\_5\_on\_a\_sphere.ipynb}. Accessed 12 June 2026:
\url{https://google-deepmind.github.io/alphaevolve_repository_of_problems/problems/60.html}.

\bibitem{AlphaEvolve}
A.~Novikov, N.~V\~u, M.~Eisenberger, E.~Dupont, P.-S. Huang, A.~Z. Wagner,
S.~Shirobokov, B.~Kozlovskii, F.~J.~R. Ruiz, A.~Mehrabian, M.~P. Kumar,
A.~See, S.~Chaudhuri, G.~Holland, A.~Davies, S.~Nowozin, P.~Kohli, and
M.~Balog,
\emph{AlphaEvolve: a coding agent for scientific and algorithmic discovery},
arXiv:2506.13131 (2025).

\bibitem{SukWhite2025}
A.~Suk and E.~P. White,
\emph{A note on the no-$(d+2)$-on-a-sphere problem},
in: 41st International Symposium on Computational Geometry (SoCG 2025),
LIPIcs \textbf{332} (2025), 76:1--76:8; arXiv:2412.02866.

\bibitem{Szabo2025}
D.~R. Szab\'o,
\emph{Rational normal curves as no-$(d+2)$-on-$Q$-quadric sets},
arXiv:2511.03526 (2025).

\bibitem{Thiele1995}
T.~Thiele,
\emph{Geometric Selection Problems and Hypergraphs},
Dissertation, Freie Universit\"at Berlin, 1995.

\end{thebibliography}
\endgroup

\end{document}
